\begin{proposition}[topologically trivial twists pull back trivially]
\hypertarget{tag-03IO}{}\label{wbII:prop:pic0-pullback}
\StatementRecord{wbII:prop:pic0-pullback}{proved}{03IO}{the definitions and constructions stated in the proof}

\ Assume that $\omega_X$ is trivial on a neighbourhood of $W$.
For every $L\in\operatorname{Pic}^0(X)$, put
\[
 A_L=(L^\vee\otimes\omega_X)|_W.
\]
Then
\[
 \eta^*A_L\simeq\mathcal O_{\mathrm{dP}_6}.
\tag{W.18}\hypertarget{tag-03FV}{}\label{wbII:eq:W:5.17a}
\]
No assertion about $H^2(X;\mathbb Z)$ is used.
\end{proposition}



\begin{proof}
\begin{enumerate}[label={$\langle1\rangle$\arabic*.},leftmargin=*]
\item 
\noindent\textbf{Step 1.} Every line bundle in $\operatorname{Pic}^0(X)$ has $c_1(L)=0$.  Naturality
therefore gives $c_1(\eta^*i^*L)=0$.  The degree-six del Pezzo surface is the
blow-up of $\mathbb P^2$ in three noncollinear points, so
$H^1(\mathrm{dP}_6,\mathcal O)=0$.  The exponential sequence consequently
makes
\[
 c_1:\operatorname{Pic}(\mathrm{dP}_6)\longrightarrow
 H^2(\mathrm{dP}_6;\mathbb Z)
\]
injective.  Hence $\eta^*i^*L$ is trivial.  The assumed neighbourhood
trivialization of $\omega_X$ also pulls back trivially, which proves
\eqref{wbII:eq:W:5.17a}.

\end{enumerate}
\end{proof}